% Quiz 05 --- answer key with marking notes. DO NOT HAND OUT.
%
% Question 1's numbers are measured: Leiden (igraph, modularity objective,
% best of three runs) on Erdos--Renyi networks, N = 10,000, three draws per
% <k>. Question 2's result was checked numerically on a simulated
% degree-corrected network (N = 300, K = 3): the simplified log-likelihood,
% the estimates and the final formula agree with the Poisson likelihood
% computed pair by pair.
%
% Compile with: xelatex -interaction=nonstopmode solutions.tex   (run twice)
\documentclass[a4paper, 11pt]{extarticle}
% Shared preamble for Quiz 05 (quiz05.tex and solutions.tex).
% Compile with: xelatex -interaction=nonstopmode <file>.tex   (run twice)
%
% Same house look as the pen-and-paper sheets in
% lecture-note/*/pen-and-paper/: Charter, A4, generous leading, room to draw.

\usepackage[top=0.7in, bottom=0.6in, left=0.8in, right=0.8in]{geometry}
\usepackage{amsmath}
\usepackage{amssymb}
\usepackage{array}
\usepackage{graphicx}
\usepackage[hidelinks]{hyperref}
\usepackage{tcolorbox}
\usepackage{enumitem}
\usepackage{etoolbox}
\usepackage{tikz}
\usepackage{fontspec}
\usetikzlibrary{calc,positioning,arrows.meta}
\definecolor{pltblue}{HTML}{1F77B4}
\definecolor{rulegrey}{HTML}{6A6D75}

% Body face. Charter, and the four ways to get it. XCharter *is* Charter and
% ships with a full TeX Live, so it is asked for first and by file: a CI runner
% answers yes to \IfFontExistsTF{Charter}, then cannot typeset it.
\IfFontExistsTF{XCharter-Roman.otf}{
  \setmainfont{XCharter}[
    Extension      = .otf,
    UprightFont    = *-Roman,
    BoldFont       = *-Bold,
    ItalicFont     = *-Italic,
    BoldItalicFont = *-BoldItalic]
}{
  \IfFontExistsTF{Charter}{\setmainfont{Charter}}{
    \IfFontExistsTF{Iowan Old Style}{\setmainfont{Iowan Old Style}}{
      \IfFontExistsTF{Georgia}{\setmainfont{Georgia}}{}}}}

\setlength{\parindent}{0pt}
\setlength{\parskip}{0.4em}
\setlength{\emergencystretch}{1.2em}
\raggedbottom
\pagestyle{empty}

% ---------------------------------------------------------------------------
% The submission link.
%
% Both the QR square and the printed address are the course's own short link,
% not the form's: paper outlives any one form. Caddy on the course droplet
% points go.skojaku.com/ans-quiz05 at the current Google Form
% (/etc/caddy/Caddyfile, block go.skojaku.com). Swapping the form is a line on
% the server, not a reprint.
%
% If the link ever changes, rebuild the square with
%
%   uvx --from segno segno --output=quiz05-form-qr.png --scale=20 --border=1 \
%       --error=m "<the url>"
%
% PNG rather than PDF for the square: segno's PDF is too spare for xdvipdfmx to
% include, and a QR is squares, so 700px across a 3cm box is already sharper
% than the printer.
% ---------------------------------------------------------------------------
\newcommand{\formlink}{https://go.skojaku.com/ans-quiz05}
\newcommand{\formshort}{\href{\formlink}{\texttt{go.skojaku.com/ans-quiz05}}}

% Quiz header. #1 = title, #2 = the line under it, #3 = term.
\newcommand{\quizhead}[3]{%
  {\footnotesize\textsc{Advanced Network Science} \hfill \textsc{#3}}\par
  \vspace{-0.2em}\rule{\textwidth}{1pt}\par
  \vspace{0.3em}
  {\LARGE\bfseries #1}\par
  \vspace{0.1em}
  {\small #2}\par
  \vspace{0.7em}
  {\small Name \underline{\hspace{6cm}} \hfill Date \underline{\hspace{3.2cm}}}\par
  \vspace{0.7em}}

% A question heading with its point value.
\newcommand{\question}[3]{%
  \vspace{0.5em}
  {\large\bfseries Question #1 \quad\normalsize\mdseries (#2 points)}\par
  \vspace{-0.3em}{\color{rulegrey}\rule{\textwidth}{0.6pt}}\par
  \vspace{0.1em}
  {\bfseries #3}\par
  \vspace{0.2em}}

% An empty box to draw or write in. #1 = height, #2 = caption under the rule.
\newcommand{\drawbox}[2]{%
  \begin{tcolorbox}[colback=white, colframe=rulegrey, boxrule=0.5pt,
                    sharp corners, height=#1, valign=top,
                    boxsep=2pt, left=4pt, right=4pt, top=3pt, bottom=3pt]
  {\footnotesize\color{rulegrey} #2}
  \end{tcolorbox}}

% Ruled writing lines, #1 of them.
\newcommand{\writelines}[1]{%
  \vspace{0.3em}
  \foreach \i in {1,...,#1}{%
    {\color{rulegrey}\rule{\textwidth}{0.4pt}}\par\vspace{1.5em}}}

\newgeometry{top=0.4in, bottom=0.4in, left=0.75in, right=0.75in}

\renewcommand{\question}[3]{\vspace{0.3em}
  {\large\bfseries Question #1 \quad\normalsize\mdseries (#2 points) --- #3}\par
  \vspace{-0.3em}{\color{rulegrey}\rule{\textwidth}{0.6pt}}\par\vspace{0.1em}}
\newcommand{\ans}[1]{{\color{pltblue}\bfseries #1}}
\newcommand{\inr}{\sum_{i:\,c_i = r}}

\begin{document}

\quizhead{Quiz 5 \; --- \; Answer key}
        {Instructor copy \quad$\cdot$\quad do not hand out}
        {Fall 2026}

\question{1}{4}{Communities in pure noise.}

Measured, Leiden on Erd\H{o}s--R\'enyi networks with $N = 10{,}000$:

\begin{center}
{\renewcommand{\arraystretch}{1.1}
\begin{tabular}{|c||*{8}{c|}}
\hline
$\langle k \rangle$ & 1 & 2 & 4 & 8 & 16 & 32 & 64 & 128 \\ \hline
$Q_{\max}$ & 0.99 & 0.85 & 0.56 & 0.36 & 0.24 & 0.17 & 0.11 & 0.08 \\ \hline
\end{tabular}}
\end{center}

(a) [2] $Q_{\max}$ \ans{goes down}, steadily, as $\langle k \rangle$ grows.
Very sparse: \ans{toward 1}. The network falls apart into small trees, and
each tree is a community with no edge leaving it. Complete:
\ans{exactly 0}. Every group of a complete graph has exactly the edges chance
predicts or fewer, so the best partition is everyone in one group, $Q = 0$.
Marking: the direction, with a falling sketch, $+1$; the two ends, $+0.5$ each.
A sketch that falls but is not labelled at the ends keeps the first point only.

(b) [2] \emph{Why.} Modularity rewards groups that hold more edges than chance
predicts. In a sparse network chance is lumpy: a node has a handful of edges,
the network is tree-like, and it can be cut into chunks with almost no edges
between them. As $\langle k \rangle$ grows, each node's edges spread evenly
over any group you propose, the fluctuations average out (relative size
$\sim 1/\sqrt{\text{expected count}}$), and no group beats chance by much.
\emph{The check.} Compare against randomized networks of the \ans{same size and
density} (or the same degree sequence): pure noise at $N = 10{,}000$,
$\langle k \rangle = 4$ already gives $Q_{\max} \approx 0.56$, so $0.45$ is no
evidence of communities at all. Marking: the reason $+1$; the check $+1$.

``Sparse networks have fewer edges, so less to cut'' with no mention of chance
or fluctuation: $+0.5$ for the reason. ``$Q$ cannot be compared across
networks'' with no baseline named: $+0.5$ for the check. The student does not
need the number 0.56; naming the right comparison is the point.

\question{2}{6}{The likelihood of the degree-corrected block model.}

(a) [2] Poisson pmf, summed over pairs, with the sheet's convention and the
$\log A_{ij}!$ term dropped (it depends on none of $c, \theta, w$):
\[
  \mathcal{L}
  = \frac12 \sum_i \sum_j \Big[ A_{ij} \log\!\big(\theta_i\theta_j w_{c_ic_j}\big)
      - \theta_i\theta_j w_{c_ic_j} \Big] .
\]
Three pieces. $\tfrac12\sum_{ij} A_{ij}(\log\theta_i + \log\theta_j) =
\sum_i k_i \log\theta_i$ (Hint 2, and $A_{ij} = A_{ji}$).
$\tfrac12\sum_{ij} A_{ij} \log w_{c_ic_j} = \tfrac12 \sum_{rs} m_{rs}\log w_{rs}$
(Hint 3). $\tfrac12\sum_{ij}\theta_i\theta_j w_{c_ic_j}
= \tfrac12\sum_{rs} w_{rs} \big(\inr \theta_i\big)\big(\sum_{j:\,c_j=s}\theta_j\big)
= \tfrac12\sum_{rs} w_{rs}$ (the normalization). So
\[
  \mathcal{L}(c, \theta, w) = {\color{pltblue} \sum_i k_i \log\theta_i
  + \frac12 \sum_{r,s} \big( m_{rs} \log w_{rs} - w_{rs} \big)} .
\]
Marking: the log-likelihood as a sum over pairs, $+1$ (keeping $\log A_{ij}!$
or writing $\sum_{i<j}$ is fine); the simplified form, $+1$, half for the
$\theta$ term and half for the $w$ terms with the last one collapsed.

(b) [2] $\partial\mathcal{L}/\partial w_{rs} = m_{rs}/w_{rs} - 1 = 0$, so
\ans{$\hat w_{rs} = m_{rs}$}. (For $r \ne s$, $w_{rs}$ and $w_{sr}$ are one
parameter that appears twice; the factor 2 cancels.) With a multiplier per
group, $k_i/\theta_i - \mu_{c_i} = 0$, so $\theta_i = k_i/\mu_r$; summing over
group $r$ gives $\mu_r = \kappa_r$, so
\ans{$\hat\theta_i = k_i / \kappa_{c_i}$}: node $i$'s share of its group's
degree. Check: node $i$'s expected degree is
$\hat\theta_i \sum_s \hat w_{c_i s} = (k_i/\kappa_r)\,\kappa_r = k_i$; the
model reproduces every degree.
Marking: $+1$ each.

\emph{Other route, full credit.} Differentiating the unsimplified form without
the constraint gives $k_i/\theta_i = \sum_j \theta_j w_{c_ic_j}$ and
$w_{rs} = m_{rs}/(\Theta_r\Theta_s)$, $\Theta_r = \inr\theta_i$, hence
$\theta_i = k_i\Theta_r/\kappa_r$: a family of equally good solutions, one
per scale. Imposing $\Theta_r = 1$ picks $k_i/\kappa_r$.
\emph{Slips.} $\hat\theta_i = k_i/(\text{anything not } \kappa_{c_i})$ with the
normalization ignored: $+0.5$. $\hat w_{rs} = m_{rs}/(n_r n_s)$, the plain
SBM's answer: $0$ for that point.

(c) [2] Substitute:
\[
  \mathcal{L}
  = \sum_i k_i \log k_i - \sum_r \kappa_r \log\kappa_r
    + \frac12 \sum_{r,s} m_{rs}\log m_{rs} - M .
\]
$\sum_i k_i \log k_i$ and $M$ do not depend on $c$. Since
$\sum_s m_{rs} = \kappa_r$,
$\sum_r \kappa_r\log\kappa_r = \tfrac12\sum_{rs} m_{rs}(\log\kappa_r + \log\kappa_s)$, and
\[
  \mathcal{L}(c) = {\color{pltblue} \frac12 \sum_{r,s} m_{rs}
  \log \frac{m_{rs}}{\kappa_r \kappa_s}} + \text{const.}
\]
Marking: the substitution, correct in any form, $+1$; the constants named and
dropped, leaving $c$ only through $m_{rs}$ and $\kappa_r$, $+1$. The compact
form is not required:
$\tfrac12\sum_{rs} m_{rs}\log m_{rs} - \sum_r \kappa_r\log\kappa_r$ is full
marks. With or without the $\tfrac12$ is full marks (Karrer and Newman drop it).

\emph{Done early (not graded).} The plain SBM has $n_r n_s$ where this has
$\kappa_r\kappa_s$: a block is judged dense against the degrees of its members,
not their number. A group of hubs has many edges, but $\kappa_r$ is large too,
so it no longer looks surprisingly dense; the groups are free to follow who
connects to whom.

\end{document}
